%%%%%%%%%%%%%%%%%%%%%%%%%%%%%%%%%%%%%%%%%
%  Iscp Handbook EE department 
% The original template (the Legrand Orange Book Template) can be found here --> http://www.latextemplates.com/template/the-legrand-orange-book
%
% Original author of the Legrand Orange Book Template:
% Mathias Legrand (legrand.mathias@gmail.com) with modifications by:
% Vel (vel@latextemplates.com)
%
% Original License:
% CC BY-NC-SA 3.0 (http://creativecommons.org/licenses/by-nc-sa/3.0/)
% just a mere adaptation of https://www.overleaf.com/articles/clustering-the-interstellar-medium/mtthgyyfrdkn
%%%%%%%%%%%%%%%%%%%%%%%%%%%%%%%%%%%%%%%%%
 
%----------------------------------------------------------------------------------------
%	PACKAGES AND OTHER DOCUMENT CONFIGURATIONS
%----------------------------------------------------------------------------------------

\documentclass[openany]{book} % Default font size and left-justified equations

\usepackage[top=3cm,bottom=3cm,left=3.2cm,right=3.2cm,headsep=10pt,letterpaper]{geometry} % Page margins

\usepackage{xcolor} % Required for specifying colors by name
\definecolor{ocre}{RGB}{52,177,201} % Define the orange color used for highlighting throughout the book

% Font Settings
\usepackage{avant} % Use the Avantgarde font for headings
%\usepackage{times} % Use the Times font for headings
\usepackage{mathptmx} % Use the Adobe Times Roman as the default text font together with math symbols from the Sym­bol, Chancery and Com­puter Modern fonts

\usepackage{microtype} % Slightly tweak font spacing for aesthetics
\usepackage[utf8]{inputenc} % Required for including letters with accents
\usepackage[T1]{fontenc} % Use 8-bit encoding that has 256 glyphs

\usepackage[object=vectorian]{pgfornament} %for ornamental lines 
\newcommand{\sectionlinetwo}[2]{%
  \nointerlineskip \vspace{.5\baselineskip}\hspace{\fill}
  {\resizebox{0.5\linewidth}{1.2ex}
    {\pgfornament[color = #1]{#2}
    }}%
    \hspace{\fill}
    \par\nointerlineskip \vspace{.5\baselineskip}
  }
\usepackage{tikz}
\newcommand{\photo}[3]{%
	\tikz\node[circle,draw,inner sep=#1,text=white,path picture={\node at (path picture bounding box.center){\includegraphics[width=#2]{#3}};}]{};
}%

% Bibliography
\usepackage[style=alphabetic,sorting=nyt,sortcites=true,autopunct=true,babel=hyphen,hyperref=true,abbreviate=false,backref=true,backend=biber]{biblatex}


\usepackage[colorlinks = true,
	    linkcolor = blue,
	    urlcolor  = blue,
	    citecolor = blue,
	    anchorcolor = blue]{hyperref}
\usepackage{miama}
\usepackage[T1]{fontenc}

\usepackage{csvsimple} % for reading data directly from csv
\begin{filecontents*}{./iscp/test.csv}
imgA,nameA,emailA,imgB,nameB,emailB,imgB,nameC,emailC,imgD,nameD,emailD
\end{filecontents*}

\addbibresource{bibliography.bib} % BibTeX bibliography file
\defbibheading{bibempty}{}

%%%%%%%%%%%%%%%%%%%%%%%%%%%%%%%%%%%%%%%%%
% This is based on the Legrand Orange Book
% Structural Definitions File
%
% The original template (the Legrand Orange Book Template) can be found here --> http://www.latextemplates.com/template/the-legrand-orange-book
%
% Original author of the Legrand Orange Book Template::
% Mathias Legrand (legrand.mathias@gmail.com) with modifications by:
% Vel (vel@latextemplates.com)
%
% Original License:
% CC BY-NC-SA 3.0 (http://creativecommons.org/licenses/by-nc-sa/3.0/)
%
%%%%%%%%%%%%%%%%%%%%%%%%%%%%%%%%%%%%%%%%%
%----------------------------------------------------------------------------------------
%	VARIOUS REQUIRED PACKAGES
%----------------------------------------------------------------------------------------

\usepackage{titlesec} % Allows customization of titles

\usepackage{graphicx} % Required for including pictures
\graphicspath{{Pictures/}} % Specifies the directory where pictures are stored

\usepackage{lipsum} % Inserts dummy text

\usepackage{tikz} % Required for drawing custom shapes

\usepackage[english]{babel} % English language/hyphenation

\usepackage{enumitem} % Customize lists
\setlist{nolistsep} % Reduce spacing between bullet points and numbered lists

\usepackage{booktabs} % Required for nicer horizontal rules in tables

\usepackage{eso-pic} % Required for specifying an image background in the title page

%----------------------------------------------------------------------------------------
%	MAIN TABLE OF CONTENTS
%----------------------------------------------------------------------------------------

\usepackage{titletoc} % Required for manipulating the table of contents

\contentsmargin{0cm} % Removes the default margin
% Chapter text styling
\titlecontents{chapter}[1.25cm] % Indentation
{\addvspace{15pt}\large\sffamily\bfseries} % Spacing and font options for chapters
{\color{ocre!60}\contentslabel[\Large\thecontentslabel]{1.25cm}\color{ocre}} % Chapter number
{}  
{\color{ocre!60}\normalsize\sffamily\bfseries\;\titlerule*[.5pc]{.}\;\thecontentspage} % Page number
% Section text styling
\titlecontents{section}[1.25cm] % Indentation
{\addvspace{5pt}\sffamily\bfseries} % Spacing and font options for sections
{\contentslabel[\thecontentslabel]{1.25cm}} % Section number
{}
{\sffamily\hfill\color{black}\thecontentspage} % Page number
[]
% Subsection text styling
\titlecontents{subsection}[1.25cm] % Indentation
{\addvspace{1pt}\sffamily\small} % Spacing and font options for subsections
{\contentslabel[\thecontentslabel]{1.25cm}} % Subsection number
{}
{\sffamily\;\titlerule*[.5pc]{.}\;\thecontentspage} % Page number
[] 

%----------------------------------------------------------------------------------------
%	MINI TABLE OF CONTENTS IN CHAPTER HEADS
%----------------------------------------------------------------------------------------

% Section text styling
\titlecontents{lsection}[0em] % Indendating
{\footnotesize\sffamily} % Font settings
{}
{}
{}

% Subsection text styling
\titlecontents{lsubsection}[.5em] % Indentation
{\normalfont\footnotesize\sffamily} % Font settings
{}
{}
{}
 
%----------------------------------------------------------------------------------------
%	PAGE HEADERS
%----------------------------------------------------------------------------------------

\usepackage{fancyhdr} % Required for header and footer configuration

\pagestyle{fancy}
\renewcommand{\chaptermark}[1]{\markboth{\sffamily\normalsize\bfseries\chaptername\ \thechapter.\ #1}{}} % Chapter text font settings
\renewcommand{\sectionmark}[1]{\markright{\sffamily\normalsize\thesection\hspace{5pt}#1}{}} % Section text font settings
\fancyhf{} \fancyhead[LE,RO]{\sffamily\normalsize\thepage} % Font setting for the page number in the header
\fancyhead[LO]{\rightmark} % Print the nearest section name on the left side of odd pages
\fancyhead[RE]{\leftmark} % Print the current chapter name on the right side of even pages
\renewcommand{\headrulewidth}{0.5pt} % Width of the rule under the header
\addtolength{\headheight}{2.5pt} % Increase the spacing around the header slightly
\renewcommand{\footrulewidth}{0pt} % Removes the rule in the footer
\fancypagestyle{plain}{\fancyhead{}\renewcommand{\headrulewidth}{0pt}} % Style for when a plain pagestyle is specified

% Removes the header from odd empty pages at the end of chapters
\makeatletter
\renewcommand{\cleardoublepage}{
\clearpage\ifodd\c@page\else
\hbox{}
\vspace*{\fill}
\thispagestyle{empty}
\newpage
\fi}

%----------------------------------------------------------------------------------------
%	THEOREM STYLES
%----------------------------------------------------------------------------------------

\usepackage{amsmath,amsfonts,amssymb,amsthm} % For math equations, theorems, symbols, etc

\newcommand{\intoo}[2]{\mathopen{]}#1\,;#2\mathclose{[}}
\newcommand{\ud}{\mathop{\mathrm{{}d}}\mathopen{}}
\newcommand{\intff}[2]{\mathopen{[}#1\,;#2\mathclose{]}}
\newtheorem{notation}{Notation}[chapter]

%%%%%%%%%%%%%%%%%%%%%%%%%%%%%%%%%%%%%%%%%%%%%%%%%%%%%%%%%%%%%%%%%%%%%%%%%%%
%%%%%%%%%%%%%%%%%%%% dedicated to boxed/framed environements %%%%%%%%%%%%%%
%%%%%%%%%%%%%%%%%%%%%%%%%%%%%%%%%%%%%%%%%%%%%%%%%%%%%%%%%%%%%%%%%%%%%%%%%%%
\newtheoremstyle{ocrenumbox}% % Theorem style name
{0pt}% Space above
{0pt}% Space below
{\normalfont}% % Body font
{}% Indent amount
{\small\bf\sffamily\color{ocre}}% % Theorem head font
{\;}% Punctuation after theorem head
{0.25em}% Space after theorem head
{\small\sffamily\color{ocre}\thmname{#1}\nobreakspace\thmnumber{\@ifnotempty{#1}{}\@upn{#2}}% Theorem text (e.g. Theorem 2.1)
\thmnote{\nobreakspace\the\thm@notefont\sffamily\bfseries\color{black}---\nobreakspace#3.}} % Optional theorem note
\renewcommand{\qedsymbol}{$\blacksquare$}% Optional qed square

\newtheoremstyle{blacknumex}% Theorem style name
{5pt}% Space above
{5pt}% Space below
{\normalfont}% Body font
{} % Indent amount
{\small\bf\sffamily}% Theorem head font
{\;}% Punctuation after theorem head
{0.25em}% Space after theorem head
{\small\sffamily{\tiny\ensuremath{\blacksquare}}\nobreakspace\thmname{#1}\nobreakspace\thmnumber{\@ifnotempty{#1}{}\@upn{#2}}% Theorem text (e.g. Theorem 2.1)
\thmnote{\nobreakspace\the\thm@notefont\sffamily\bfseries---\nobreakspace#3.}}% Optional theorem note

\newtheoremstyle{blacknumbox} % Theorem style name
{0pt}% Space above
{0pt}% Space below
{\normalfont}% Body font
{}% Indent amount
{\small\bf\sffamily}% Theorem head font
{\;}% Punctuation after theorem head
{0.25em}% Space after theorem head
{\small\sffamily\thmname{#1}\nobreakspace\thmnumber{\@ifnotempty{#1}{}\@upn{#2}}% Theorem text (e.g. Theorem 2.1)
\thmnote{\nobreakspace\the\thm@notefont\sffamily\bfseries---\nobreakspace#3.}}% Optional theorem note

%%%%%%%%%%%%%%%%%%%%%%%%%%%%%%%%%%%%%%%%%%%%%%%%%%%%%%%%%%%%%%%%%%%%%%%%%%%
%%%%%%%%%%%%% dedicated to non-boxed/non-framed environements %%%%%%%%%%%%%
%%%%%%%%%%%%%%%%%%%%%%%%%%%%%%%%%%%%%%%%%%%%%%%%%%%%%%%%%%%%%%%%%%%%%%%%%%%
\newtheoremstyle{ocrenum}% % Theorem style name
{5pt}% Space above
{5pt}% Space below
{\normalfont}% % Body font
{}% Indent amount
{\small\bf\sffamily\color{ocre}}% % Theorem head font
{\;}% Punctuation after theorem head
{0.25em}% Space after theorem head
{\small\sffamily\color{ocre}\thmname{#1}\nobreakspace\thmnumber{\@ifnotempty{#1}{}\@upn{#2}}% Theorem text (e.g. Theorem 2.1)
\thmnote{\nobreakspace\the\thm@notefont\sffamily\bfseries\color{black}---\nobreakspace#3.}} % Optional theorem note
\renewcommand{\qedsymbol}{$\blacksquare$}% Optional qed square
\makeatother

% Defines the theorem text style for each type of theorem to one of the three styles above
\newcounter{dummy} 
\numberwithin{dummy}{section}
\theoremstyle{ocrenumbox}
\newtheorem{theoremeT}[dummy]{Theorem}
\newtheorem{problem}{Problem}[chapter]
\newtheorem{exerciseT}{Exercise}[chapter]
\theoremstyle{blacknumex}
\newtheorem{exampleT}{Example}[chapter]
\theoremstyle{blacknumbox}
\newtheorem{vocabulary}{Vocabulary}[chapter]
\newtheorem{definitionT}{Definition}[section]
\newtheorem{corollaryT}[dummy]{Corollary}
\theoremstyle{ocrenum}
\newtheorem{proposition}[dummy]{Proposition}

%----------------------------------------------------------------------------------------
%	DEFINITION OF COLORED BOXES
%----------------------------------------------------------------------------------------

\RequirePackage[framemethod=default]{mdframed} % Required for creating the theorem, definition, exercise and corollary boxes

% Theorem box
\newmdenv[skipabove=7pt,
skipbelow=7pt,
backgroundcolor=black!5,
linecolor=ocre,
innerleftmargin=5pt,
innerrightmargin=5pt,
innertopmargin=5pt,
leftmargin=0cm,
rightmargin=0cm,
innerbottommargin=5pt]{tBox}

% Exercise box	  
\newmdenv[skipabove=7pt,
skipbelow=7pt,
rightline=false,
leftline=true,
topline=false,
bottomline=false,
backgroundcolor=ocre!10,
linecolor=ocre,
innerleftmargin=5pt,
innerrightmargin=5pt,
innertopmargin=5pt,
innerbottommargin=5pt,
leftmargin=0cm,
rightmargin=0cm,
linewidth=4pt]{eBox}	

% Definition box
\newmdenv[skipabove=7pt,
skipbelow=7pt,
rightline=false,
leftline=true,
topline=false,
bottomline=false,
linecolor=ocre,
innerleftmargin=5pt,
innerrightmargin=5pt,
innertopmargin=0pt,
leftmargin=0cm,
rightmargin=0cm,
linewidth=4pt,
innerbottommargin=0pt]{dBox}	

% Corollary box
\newmdenv[skipabove=7pt,
skipbelow=7pt,
rightline=false,
leftline=true,
topline=false,
bottomline=false,
linecolor=gray,
backgroundcolor=black!5,
innerleftmargin=5pt,
innerrightmargin=5pt,
innertopmargin=5pt,
leftmargin=0cm,
rightmargin=0cm,
linewidth=4pt,
innerbottommargin=5pt]{cBox}

% Creates an environment for each type of theorem and assigns it a theorem text style from the "Theorem Styles" section above and a colored box from above
\newenvironment{theorem}{\begin{tBox}\begin{theoremeT}}{\end{theoremeT}\end{tBox}}
\newenvironment{exercise}{\begin{eBox}\begin{exerciseT}}{\hfill{\color{ocre}\tiny\ensuremath{\blacksquare}}\end{exerciseT}\end{eBox}}				  
\newenvironment{definition}{\begin{dBox}\begin{definitionT}}{\end{definitionT}\end{dBox}}	
\newenvironment{example}{\begin{exampleT}}{\hfill{\tiny\ensuremath{\blacksquare}}\end{exampleT}}		
\newenvironment{corollary}{\begin{cBox}\begin{corollaryT}}{\end{corollaryT}\end{cBox}}	

%----------------------------------------------------------------------------------------
%	REMARK ENVIRONMENT
%----------------------------------------------------------------------------------------

\newenvironment{remark}{\par\vspace{10pt}\small % Vertical white space above the remark and smaller font size
\begin{list}{}{
\leftmargin=35pt % Indentation on the left
\rightmargin=25pt}\item\ignorespaces % Indentation on the right
\makebox[-2.5pt]{\begin{tikzpicture}[overlay]
\node[draw=ocre!60,line width=1pt,circle,fill=ocre!25,font=\sffamily\bfseries,inner sep=2pt,outer sep=0pt] at (-15pt,0pt){\textcolor{ocre}{R}};\end{tikzpicture}} % Orange R in a circle
\advance\baselineskip -1pt}{\end{list}\vskip5pt} % Tighter line spacing and white space after remark

%----------------------------------------------------------------------------------------
%	SECTION NUMBERING IN THE MARGIN
%----------------------------------------------------------------------------------------

\makeatletter
\renewcommand{\@seccntformat}[1]{\llap{\textcolor{ocre}{\csname the#1\endcsname}\hspace{1em}}}                    
\renewcommand{\section}{\@startsection{section}{1}{\z@}
{-4ex \@plus -1ex \@minus -.4ex}
{1ex \@plus.2ex }
{\normalfont\large\sffamily\bfseries}}
\renewcommand{\subsection}{\@startsection {subsection}{2}{\z@}
{-3ex \@plus -0.1ex \@minus -.4ex}
{0.5ex \@plus.2ex }
{\normalfont\sffamily\bfseries}}
\renewcommand{\subsubsection}{\@startsection {subsubsection}{3}{\z@}
{-2ex \@plus -0.1ex \@minus -.2ex}
{.2ex \@plus.2ex }
{\normalfont\small\sffamily\bfseries}}                        
\renewcommand\paragraph{\@startsection{paragraph}{4}{\z@}
{-2ex \@plus-.2ex \@minus .2ex}
{.1ex}
{\normalfont\small\sffamily\bfseries}}

%----------------------------------------------------------------------------------------
%	HYPERLINKS IN THE DOCUMENTS
%----------------------------------------------------------------------------------------

% For an unclear reason, the package should be loaded now and not later
\usepackage{hyperref}
\hypersetup{hidelinks,backref=true,pagebackref=true,hyperindex=true,colorlinks=false,breaklinks=true,urlcolor= ocre,bookmarks=true,bookmarksopen=false,pdftitle={Title},pdfauthor={Author}}

%----------------------------------------------------------------------------------------
%	CHAPTER HEADINGS
%----------------------------------------------------------------------------------------

% The set-up below should be (sadly) manually adapted to the overall margin page septup controlled by the geometry package loaded in the main.tex document. It is possible to implement below the dimensions used in the goemetry package (top,bottom,left,right)... TO BE DONE

\newcommand{\thechapterimage}{}
\newcommand{\chapterimage}[1]{\renewcommand{\thechapterimage}{#1}}

% Numbered chapters with mini tableofcontents
\def\thechapter{\arabic{chapter}}
\def\@makechapterhead#1{
\thispagestyle{empty}
{\centering \normalfont\sffamily
\ifnum \c@secnumdepth >\m@ne
\if@mainmatter
\startcontents
\begin{tikzpicture}[remember picture,overlay]
\node at (current page.north west)
{\begin{tikzpicture}[remember picture,overlay]
\node[anchor=north west,inner sep=0pt] at (0,0) {\includegraphics[width=\paperwidth]{\thechapterimage}};
%%%%%%%%%%%%%%%%%%%%%%%%%%%%%%%%%%%%%%%%%%%%%%%%%%%%%%%%%%%%%%%%%%%%%%%%%%%%%%%%%%%%%
% Commenting the 3 lines below removes the small contents box in the chapter heading
%\fill[color=ocre!10!white,opacity=.6] (1cm,0) rectangle (8cm,-7cm);
%\node[anchor=north west] at (1.1cm,.35cm) {\parbox[t][8cm][t]{6.5cm}{\huge\bfseries\flushleft \printcontents{l}{1}{\setcounter{tocdepth}{2}}}};
\draw[anchor=west] (5cm,-9cm) node [rounded corners=20pt,fill=ocre!10!white,text opacity=1,draw=ocre,draw opacity=1,line width=1.5pt,fill opacity=.6,inner sep=12pt]{\huge\sffamily\bfseries\textcolor{black}{\thechapter. #1\strut\makebox[22cm]{}}};
%%%%%%%%%%%%%%%%%%%%%%%%%%%%%%%%%%%%%%%%%%%%%%%%%%%%%%%%%%%%%%%%%%%%%%%%%%%%%%%%%%%%%
\end{tikzpicture}};
\end{tikzpicture}}
\par\vspace*{230\p@}
\fi
\fi}

% Unnumbered chapters without mini tableofcontents (could be added though) 
\def\@makeschapterhead#1{
\thispagestyle{empty}
{\centering \normalfont\sffamily
\ifnum \c@secnumdepth >\m@ne
\if@mainmatter
\begin{tikzpicture}[remember picture,overlay]
\node at (current page.north west)
{\begin{tikzpicture}[remember picture,overlay]
\node[anchor=north west,inner sep=0pt] at (0,0) {\includegraphics[width=\paperwidth]{\thechapterimage}};
\draw[anchor=west] (5cm,-9cm) node [rounded corners=20pt,fill=ocre!10!white,fill opacity=.6,inner sep=12pt,text opacity=1,draw=ocre,draw opacity=1,line width=1.5pt]{\huge\sffamily\bfseries\textcolor{black}{#1\strut\makebox[22cm]{}}};
\end{tikzpicture}};
\end{tikzpicture}}
\par\vspace*{230\p@}
\fi
\fi
}
\makeatother % Insert the commands.tex file which contains the majority of the structure behind the template

\begin{document}
\title{EE Handbook ISCP}

%----------------------------------------------------------------------------------------
%	TITLE PAGE
%----------------------------------------------------------------------------------------

\begingroup
\thispagestyle{empty}
\AddToShipoutPicture*{\put(0,0){\includegraphics[scale=0.42]{./pictures/iitbombayTwire.jpg}}} % Image background
\centering
\vspace*{5cm}
\par\normalfont\fontsize{35}{35}\sffamily\selectfont
\textbf{Department Handbook}\\
{\LARGE Electrical Engineering}\par % Book title
{\Huge 2019-2020 }\par % Author name
\endgroup

%----------------------------------------------------------------------------------------
%	DISCLAIMER PAGE
%----------------------------------------------------------------------------------------

\newpage
~\vfill
\thispagestyle{empty}

%\noindent Copyright \copyright\ 2014 Andrea Hidalgo\\ % Copyright notice

\noindent \textsc{DISCLAIMER}\\

%\noindent \textsc{github.com/LaurethTeX/Clustering}\\ % URL

\noindent The Institute Student Companion Program (ISCP) has acquired and presented the data in this
handbook on a best effort basis. However the correctness of the information is not guaranteed.
ISCP will not be held responsible for any inaccuracies in the document.\\

\noindent \textit{First release, June 2019} % Printing/edition date

%----------------------------------------------------------------------------------------
%	TABLE OF CONTENTS
%----------------------------------------------------------------------------------------

\chapterimage{./pictures/index_hand.jpg} % Table of contents heading image

\pagestyle{empty} % No headers

\tableofcontents % Print the table of contents itself

%\cleardoublepage % Forces the first chapter to start on an odd page so it's on the right

\pagestyle{fancy} % Print headers again

%----------------------------------------------------------------------------------------
%	CHAPTER 1
%----------------------------------------------------------------------------------------

\chapterimage{./pictures/convo.jpg} % Chapter heading image

\chapter{About The Department}


\section*{}
The Department of \href{https://www.ee.iitb.ac.in/web}{\color{blue}{Electrical Engineering (EE)}} is one of the largest departments of IIT Bombay since its inception in 1958. The department has different academic programs with about 570 undergraduate and 730 postgraduate students. The department is equipped with the state of the art experimental and computational facilities for undertaking R $\&$ D and consultancy activities in various fields.\\
\\
The EE department has a vibrant postgraduate program with strong focus on research and development. The number of postgraduate students in EE is more than that of the undergraduate students and the department attaches a lot of importance to its Masters' students as they constitute the backbone of research and development.
\\
\\The department offers M.Tech in six research areas / specializations:
\begin{itemize}
    \item Communications Engineering (EE1)
    \item Control and Computing (EE2)
    \item Power Electronics and Power Systems (EE3)
    \item Electronic Systems (EE5)
    \item Integrated Circuit and Systems (EE6)
    \item Solid State Devices (EE7)
\end{itemize}
\medskip
Faculty members of the department are recipients of many distinguished awards like Shanti Swarup Bhatnagar Prize, Prof. K. Sreenivasan Memorial Award, Prof. SVC Aiya Memorial Award, Dr. Vikram Sarabhai Research Award, Ram Lal Wadhwa Award, INAE Young Engineer Award, Alexander von Humboldt Fellowship and many others.\\
\\
Many faculty members are Editors of IEEE and other national and international journals. They are also Fellows of organizations like IEEE, IETE, INAE, IASc, NASI and INSA.
%This statement requires citation \cite{book_key}; this one is more specific \cite[122]{article_key}.
%----------------------------------------------------------------------------------------
%	CHAPTER 2
%----------------------------------------------------------------------------------------
\chapterimage{./pictures/bgf4.png}

\chapter{Message from H.O.D}

\section*{}
Congratulations on your selection for the M.Tech/Ph.D. program in EE at
IIT Bombay. As you know, the competition was very stiff and you are among
the very few students who made it. We, the faculty members, staff and
students extend a warm welcome to you.
%\begin{figure}[h]
    %\centering
    %\includegraphics[width=0.77\textwidth]{ha-gray-conv-crp.jpg}
    %\caption{Picture of the M83 galaxy, image taken from the WFC3 ERS M83 Data Products, http://archive.stsci.edu/prepds/wfc3ers/m83datalist.html}
    %\label{fig:awesome_image}
%\end{figure}
\medskip
\newline
As you are aware, ours is among the largest Electrical Engineering
departments in the country with 66 faculty members and 1360 students, of
which more than 55\% are graduate student. We have a strong academic and
research culture. We have state-of-the-art research laboratories in almost
all areas of electrical engineering and a few centres of excellence. I am
sure you will find this place academically rewarding.
\medskip
\newline
This department has a lot to give –just how much you take depends on one
person-and that is you. You will face diverse temptations but you need to
stay focused to achieve your goals. Do good work-you grow, and the
department does too. It is a win-win situation. All in all, I assure you
the time spent here will be the best years
of your life!
\newline
\begin{center}
    \textbf{Feel free to contact me if you need any help!}\\

\begin{flushright}
	{\LARGE\fmmfamily
	Baylon G Fernandes\\}
	Head Of Department\\
	head@ee.iitb.ac.in\\
\end{flushright}
\end{center}
%----------------------------------------------------------------------------------------
%	CHAPTER 3
%----------------------------------------------------------------------------------------


\chapterimage{./pictures/iscp_bg.png}
\chapter{Message from ISCP}
 %{\thechapterimage}};

\section*{}

Dear New Entrants,\\
\\
We take this opportunity to welcome you to one of the most prestigious institutes of the country. We
congratulate you on having achieved this feat. With our personal experiences we can vouch that the your stay here at the campus would be exciting. From potential leopard sightings to potential bumping into movie stars all awaits you. Wonderful all night banters, amazing wing cultures and mad birthday celebrations are a few things that you will carry from here when you leave, obviously along with the degree. You will also become a part of a culture where people want to perfect their craft and thus work day in and day out at it. Hence there will be great opportunities to learn both inside as well as outside the classrooms. Thus it is a whole new cosmos to enter and with excitements it might have a few challenges too. We at \href{https://gymkhana.iitb.ac.in/~scp/scp/index.html}{\color{blue}{Institute Students Companion Program (ISCP)}}  work towards providing you with the hacks to take care of these challenges and have a happy stay with the IIT Bombay family.\\
\\
\\
The primary objective of the Companion Programme under which the ISCP team works is to build a
relationship of trust and comfort between the final year students and the incoming students of the PG
programmes. Once this is established life at campus becomes so much easier than what it would have been
without it. The knowledge and the experience that the senior batch has gained with their stay at the campus makes the surroundings so much you that the transition becomes smooth. From the lingo on campus to the terminology` in the classroom, from the grading to the syllabus, from the profile to placements, from tagda franky to bhindi rice all becomes ingrained so much as if it were you always.\\ 
\\
\\
On campus you might always be short on time as there is so much to do and when there is so much to do time
flies at sonic speeds. Managing the academics along with extra curricular activities and your social life may seem a daunting task at times. The ISCP programme thus provides you a Student Companion with whom you
can share your academic and non- academic problems. These are self-motivated volunteers who want to
genuinely help you in tough situation as a giving back act of what they received from the programme.\\ 
\\ 
\\ You can look up to the team for any initial information in things that you are venturing out at be it academics or extra-curriculars, any academic or non- academic issues that you are facing, any sort of support, any requirement that you wanna raise up as a part of the student community and last but certainly not the least just for normal interaction because that is all the programme holds at its core.\\
\\
\\Come be a part of this immense pool of wisdom and make it more happening and diverse.\\

\bigbreak
\begin{center}
Mail to: iscp@iitb.ac.in\\
\textbf{Overall Coordinators}\\
\textbf {Institute Student Companion Programme (2019-20)}\\
\medskip
\bigbreak
\begin{center}
	\begin{tabular}{ccc}
		\photo{1cm}{35mm}{./pictures/uroosa.jpeg}
		& \photo{1cm}{35mm}{./pictures/tumul.jpeg} 
		& \photo{1cm}{28mm}{./pictures/avinash.jpeg} \\
\bigbreak
		 \textbf{Uroosa Warsi}
		&\textbf{Tumul Rai}
		&\textbf{Avinash Singh}\\
		\textit{uroosawarsi134@gmail.com}
		&\textit{tumulrai91@gmail.com}
		& \textit{Avinashindolia007@gmail.com}\\
		\textbf{7835877634}
		&\textbf{7275362979}
		&\textbf{+919058076777}\\
		\end{tabular}
%Uroosa Warsi | Tumul Rai | Avinash Singh\\
%+91-9007766390 | +91-9432152174 | +91-7008955255
%\vfill
	\\
	\medskip
	\sectionlinetwo{black}{88}
	%\pgfornament[color=black]{89}
	 % \foreach \i in {1,...,99} {\expandafter\pgfornament\expandafter{\i}\ \i \\ }
 \end{center}
 \end{center}
~\vfill
%----------------------------------------------------------------------------------------
%	CHAPTER 4
%----------------------------------------------------------------------------------------

\chapterimage{./pictures/imr_bg.png} % Chapter heading image

\chapter{Message from PGAC}

\section*{}
Hello, Friends!\\
	Congratulations on being selected to be a part of IIT Bombay and a cordial welcome to this new world.\\
You must have realized by now that you are undergoing a phase of substantive transformation, and this might be daunting for some of you. The academic curriculum of this institution might seem different and perhaps new in comparison to what you were exposed to during your undergraduate education. To appease all the apprehensions that you have, the Post Graduate Academic Council (PGAC) along with the team of (ISCP) will try to address all the queries that you will be having during your entire Master’s Programme. \\
IIT Bombay is known to offer the students a very dynamic environment and a reasonable amount of freedom so that the individual can pursue their heart’s desire. Be it academics, sports, cultural, or any other activity, you will be finding myriad opportunities to build up your personality and add value to your life. I am hoping that you will be able to explore the unending map of possibilities, push your boundaries, break all the walls and bring out the best version of yourself by the time you finish your Degree Programme.\\
I wish you good luck and hope to see you around. We are looking forward to interacting with you. \\
\\
\begin{center}
	\begin{tabular}{c}
		\photo{1cm}{35mm}{./pictures/issa.jpeg} \\
		\textbf{Himanshu Bishwash}\\
		\textbf{imr@iitb.ac.in}\\
		\textit{8698965842/9867110738}\\
		\end{tabular}
\end{center}

%----------------------------------------------------------------------------------------
%	CHAPTER 5
%----------------------------------------------------------------------------------------


\chapterimage{./pictures/labs.jpg}
\chapter{LAB Facilities}
\begin{enumerate}
    \item \href{https://www.ee.iitb.ac.in/~wel_iitb/}{\color{blue}  \textbf{Wadhwani Electronics Laboratory (WEL)}}\\
    (\textit{3rd Floor , EE Building})\\
    Professor In-charge - Prof Siddharth Tallur\\
Relevant Specializations - for all specializations in EE\\
\\
The WEL houses all the major electronics hardware activities of the Electrical Department at IIT Bombay. Around 12 major laboratory course are conducted every year in WEL. Around 650 UG students and 250 PG and DD students enroll for lab courses annually. In addition to courses, WEL is also home to all major electronics projects done by students at IIT Bombay. Various Technical events are housed here, including electronics workshops,competitions, teachers' training etc.\\

\item \href{http://www.ee.iitb.ac.in/~spilab} {\color{blue}  \textbf{Signal Processing and Instrumentation Lab} }\\
    (\textit{1st Floor, EE Building})\\
    Professor In-charge - Prof. P C Pandey\\
Relevant Specializations - EE1 , EE5\\
\\
This lab focuses on research in the areas of speech signal processing, bio-medical signal processing \& instrumentation, electronic instrumentation and embedded system design. The research problems being tackled by students in the lab include 'Enhancement of electroaryngeal speech using spectral subtraction' , 'Multi-band frequency compression for hearing impaired' , 'Speech enhancement by modification of stop consonant landmarks', 'Diagnostic information from impedance cardiograph' etc.\\


\item \href{https://www.ee.iitb.ac.in/web/research/labs/isl} {\color{blue} \textbf{Integrated Systems Laboratory}}\\
    (\textit{1st Floor, Electrical Annex Building, Opposite to EE main building})\\
    Professor In-charge - Prof. Jayanta Mukherjee, Prof. Maryam Shojaei Baghini\\
Relevant Specializations - EE1 , EE5 , EE6 , EE7\\
\\
Integrated Systems Laboratory at IIT Bombay is a simulations and testing laboratory located in the first floor of the EE Annex Building. Embedded system solutions are developed here. Primarily design and test of passive and active RF and circuits is done. EM design software like CST Microwave Studio, HFSS and Key sight ADS are installed in lab. Several Core i7 PC's with augmented RAM's to enable high end computing are also present. GPU based high end workstations for enabling fast EM solvers for design at millimetre wave frequencies are available in lab. This lab work closely with the VLSI lab for RFIC simulations.\\


\item \href{https://www.ee.iitb.ac.in/web/research/labs/emsys} {\color{blue} \textbf{Embedded System lab}}\\
    (\textit{5th floor, Girish Gaitonde (GG) building})\\
Relevant Specializations - EE5 , EE6 , EE7\\
\\
The prime focus areas of lab are system design, prototyping and evaluation starting from sensor/transducer interfacing to full system development and network of sensor nodes, signal processing system design, analysis and implementation for various sensing-based applications, AI for ASIC design, sensor nodes and networks, deep learning for signal and data Analysis, hardware-accelerated simulation. Research work is also done in agricultural, bio-medical domains.\\

\item \href{http://www.ee.iitb.ac.in/~stallur/index.php/} {\color{blue} \textbf{Applied Integrated Micro Systems (AIMS) Laboratory}}\\
    (\textit{1st Floor, Electrical Annex Building, Opposite to EE main building})\\
    Professor In-charge - Prof. Siddharth Tallur\\
Relevant Specializations - EE5 , EE6 , EE7\\
\\
AIMS lab works on  innovative instrumentation for impactful measurements. The research areas include sensor systems, hybrid integrated microsystems,  studying their underlying physics  to leverage such platforms for high resolution sensing applications.\\

\item \href{http://www.ee.iitb.ac.in/~kasturis/index.php} {\color{blue} \textbf{Photonics and Quantum Enabled Sensing Technology (P-Quest) Laboratory}}\\
    (\textit{2nd Floor, EE Building})\\
    Professor In-charge - Prof. Kasturi Saha\\
Relevant Specializations - EE1 , EE7\\
\\
P-Quest lab works on exploring precision metrology and sensing using novel interdisciplinary research in fields like nano-photonics, classical and quantum information processing and life sciences, to develop practical quantum devices via design and experimentation, thus connecting quantum theory to engineering applications.\\

\item \href{https://www.ee.iitb.ac.in/vlsi/} {\color{blue} \textbf{VLSI Design Lab}}\\
    (\textit{5th Floor, Girish Gaitonde (GG) Building})\\
Relevant Specializations - EE5 , EE6\\
\\
The VLSI Design Lab hosts all major VLSI CAD Vendor’s tool and their licenses .Few major tools in frequent use are by Synopsys, Cadence, Mentor, Agilent, Magma, Xilinx etc. The main research focus is in the area of analog and digital design. In addition to courses, this lab also hosts accounts for different courses which require hands on experience of tools. For different projects and their tape-out, there is availability of a high performance computational server to speed up the simulations.\\

\item \href{https://www.ee.iitb.ac.in/~spann/} {\color{blue} \textbf{Signal Processing and Artificial Neural Networks (SPANN)}}\\
    (\textit{3rd Floor, EE Building})\\
    Professor In-charge - Prof. S.N. Merchant\\
Relevant Specializations - EE1\\
\\
The major areas of research which are pursued in SPANN Lab include Wireless Communications, Sensor Networks, Image Processing and Signal Processing.\\
\\
\item \href{https://www.ee.iitb.ac.in/~infonet/} {\color{blue} \textbf{Information Networks Laboratory}}\\
    (\textit{2nd Floor, EE Building})\\
    Professor In-charge - Prof. Prasanna Chaporkar ,  Prof. Abhay Karandikar\\
Relevant Specializations - EE1\\
\\
Group members of the lab are pursuing research in the field of 4G and 5G cellular technologies, with an emphasis on inter-working with non-3GPP Wireless Local Area Networks (WLANs). Other important focus areas of the lab include-mechanisms for spectrum sharing and multi-cast for 4G networks, integrated access and backhaul systems, multi-connectivity in 5G networks and Software Defined Networking (SDN) for cellular networks and WLAN. The TTSL - IITB Center of Excellence in Telecommunication (TICET) facility is also a part of this lab.\\

\item \href{http://www.ee.iitb.ac.in/~tidsplab/index.php} {\color{blue} \textbf{Texas Instruments Digital Signal Processing Lab (TIDSP)}}\\
    (\textit{3rd Floor, EE Building})\\
    Professor In-charge - Prof. V.M. Gadre\\
Relevant Specializations - EE1 , EE5\\
\\
TIDSP laboratory was set up in the EE Department to support DSP hands-on projects at the undergraduate and the postgraduate levels. DSP specific hardware and software support is provided by Texas Instruments (TI) itself.\\
\item \href{https://www.ee.iitb.ac.in/~foclab/} {\color{blue} \textbf{Fiber-Optics Communication Lab}}\\
    (\textit{2nd Floor, EE Building})\\
    Professor In-charge - Prof. Kumar Appaiah , Prof. Joseph John\\
Relevant Specializations - EE1 , EE5\\
\\
This lab is dedicated to pursue research mainly in the area of optical fiber communication ( SM, MM, FM), plastic optical fiber and fiber sensing.\\

\item \href{https://www.ee.iitb.ac.in/bharticentre/index.html} {\color{blue} \textbf{Bharti Centre for Communication}}\\
    (\textit{2nd Floor, EE Building})\\
    Professor In-charge - Prof. D. Manjunath , Prof. Bikash Kumar Dey\\
Relevant Specializations - EE1\\
\\
The Bharati Centre for Communication is a centre to generate fundamental knowledge in telecommunication and allied systems. The Vision of the centre is to  be an internationally recognised contributor in moving the frontiers of knowledge through research and education, to keep technology practise in focus and to educate for innovation and leadership.\\

\item \href{http://www.ee.iitb.ac.in/~viplab/} {\color{blue} \textbf{Vision and Image Processing}}\\
    (\textit{ 1st Floor, EE Building})\\
    Professor In-charge - Prof. Rajbabu Velmurugan \\
Relevant Specializations - EE1,EE2,EE5\\
\\
This lab is dedicated to Deep Learning, Computer vision techniques. The major projects currently undertaken are related to Haptics, Biometrics, Image segmentation, super-resolution, Anomaly detection and surveilling related problems. This lab consists of more than 23 GPU, and high-performance computer to work on the mentioned techniques. Research is currently heading in the direction of surveillance such as aerial and single camera view. The task involves are human pose estimation, scene understanding, object detection, etc.\\

\item  \href{http://www.ee.iitb.ac.in/~pclab/index.php/}{\color{blue}\textbf{PC Lab}}\\
    (\textit{1st Floor, EE Building})\\
    Professor In-charge - Prof. Prasanna Chaporkar\\
Relevant Specializations - For all specializations in EE\\
PC Lab provides general computing facilities to students of Electrical Engineering and Reliability Engineering. Both Windows and Linux machines are present in the lab. In addition, a load-balanced server (Sharada) is available for heavier computational use. Software packages like MATLAB, Lyx, Scilab, Spice, Ansys, Sequel, Grace, etc. are installed on Sharada. You can use them for your (academic) work. Standard Linux/unix packages, such as LaTeX, mysql, etc are also available. 
\\


\item  \href{https://www.ee.iitb.ac.in/student/~daplab/people/index.html} {\color{blue}\textbf{Digital Audio Processing Lab}}\\
    (\textit{ 1st Floor, EE Building})\\
    Professor In-charge - Prof. Preeti Rao\\
Relevant Specializations - EE1,EE5\\
\\
This lab is based on the application of signal processing in the analysis of speech and audio. Research activities are related to spoken language assessment, music content analysis, measuring the goodness of instruments like Tabla, segmentation of instruments in the music concert and other application of speech and audio processing.\\




\item \href{https://www.ee.iitb.ac.in/student/~comlab/index.htm} {\color{blue}\textbf{Communication Lab}}\\
	(\textit{1st Floor, EE Building})\\
    Professor In-charge - Prof. Shalabh Gupta\\
Relevant Specializations - EE1 , EE5 , EE6\\
\\
Communication Lab primarily focuses upon cutting 
edge research in the area of High speed Communication 
Links. It can further be divided into different domains like  High speed Links using Optical Communication, 
Silicon Photonics, SerDes (Serialiser and Deserialiser) Links, RF Circuits and Millimetre-wave circuits and Systems.  Besides, this lab  also works in the domain of RF Electronics,  Embedded Systems and Audio Video Communication for some of its projects.\\

\item  \href{https://www.ee.iitb.ac.in/web/research/labs/networking} {\color{blue}\textbf{Networking Lab}} \\
	(\textit{2nd Floor, EE Building})\\
    Professor In-charge - Prof. D. Majunath\\
Relevant Specializations - EE1\\
\\
The work in the Networking lab deals with the theoretical aspects of queuing theory, sensor networks, applications of stochastic approximation, software routing etc. \\

\item \href{https://www.ee.iitb.ac.in/~mwave/} {\color{blue} \textbf{Microwave and Antenna Lab}}\\
    (\textit{3rd Floor, EE Building})\\
    Professor In-charge - Prof. Girish Kumar\\
Relevant Specializations - EE1\\
\\
Microwave Lab is involved in research work in the area of RF Systems, Electromagnetic Waves and Antenna Design. Primary research work is being done in different fields like Micro-strip Antenna, Microwave Integrated Circuits and Broadband Antennas.\\

\item \href{https://www.ee.iitb.ac.in/web/research/labs/ticet} {\color{blue} \textbf{TTSL-IITB Centre of Excellence in Telecommunications (TICET)}}\\
    (\textit{2nd Floor, EE Building})\\
    Professor In-charge - Prof. Abhay Karandikar\\
Relevant Specializations - EE1\\
\\
TICET focuses on state of art research in telecom relevant to Indian Service Providers in general and Tata Teleservices Limited (TTSL) in particular with special emphasis on rural wireless applications and connectivity. The research activities in this lab are related to Quality of Service and resource allocation in wired/wireless networks,TV White Space and its potential for affordable broadband access in India, Frugal 5G and rural broadband research and standardization.\\

\item \href{https://www.ee.iitb.ac.in/~asethi/research.html} {\color{blue} \textbf{Medical Deep learning and AI Lab (MeDAL)}}\\
    (\textit{1st Floor, EE Building})\\
    Professor In-charge - Prof. Amit Sethi , Prof. Manoj Gopalkrishnan\\
Relevant Specializations - EE1\\
\\
This lab is dedicated towards solving real world problems in the areas of medical imaging, radiology and pathology using deep learning architechtures. This lab houses high end computing facilities to work with large scale data (Giga pixel images) to solve various computer vision problems. Research group have colloboration with various hospitals and universities. Some ongoing engagements are with TATA memorial hospital, University of Illinois and King's college London.\\

\item \textbf{Information Systems and Radios (ISR) Lab}\\
    (\textit{2nd Floor, EE Building})\\
    Professor In-charge - Sibi Raj Pillai\\
Relevant Specializations - EE1\\
Information Systems and Radios (ISR Lab) fecilitates  research broadly in information theory and Radar Technology. Current research undertaken are wind-profiling algorithm for MST radars, GNSS Receiver Development, Information theoretical limit for Digital and Analog Hybrid communication, and Minimum energy transmission scheme for packetised transmission.
\\
\\
\item \href{https://www.ee.iitb.ac.in/student/~cadsl/} {\color{blue} \textbf{Computer Architecture And Dependable Systems Lab (CADSL)}}\\
    (\textit{2nd Floor, Electrical Engineering Annexxe})\\
    Professor In-charge - Virendra Singh\\
Relevant Specializations - EE5,EE6\\
\\
Areas of research
\begin{enumerate}
	\item Advance and futuristic architecture and system including compiler and operating system support for architecture
	\item Advanced dependable system including formal verfication and VLSI testing
	\item Computer Aided design of VLSI and hardware accelerator
\end{enumerate}
 
\item \href{https://www.ee.iitb.ac.in/~hpc/} {\color{blue} \textbf{High Performance Computing Lab (HPC)}}\\
    (\textit{4th Floor, Girish Gaitonde Building})\\
    Professor In-charge - Prof. Sachin B Patkar\\
Relevant Specializations - EE1,EE2,EE5,EE6\\
\\
Areas of research
Our research is aimed at addressing future applications and implementations of high-end   parallel and reconfigurable computing especially to Electrical Engineering related problems. Projects range from accelerating circuit simulation for digital, RF and power electronics, stereo imaging and Machine Learning on reconfigurable hardware, Network-on-Chip for Distributed Computing, and Crypt-analysis which lie at the union of mathematics, computing and Electrical Engineering. The lab caters to multiple disciplines including but not limited to Control and Computing, VLSI and Electronic Systems.

\item \href{https://www.ee.iitb.ac.in/~ccgroup} {\color{blue} \textbf{Control and Computing Laboratory}}\\
    (\textit{2nd Floor, EE Building})\\
Relevant Specializations - EE2\\
In here we pursue research in various fields. Some of them are listed below
Field of research: Differential games, Formation control and consensus in Quad-rotors, Pursuit-Evasion games, Multi-agent systems and Co-operation control, Hamiltonian systems, Stability of switched systems, Differential algebraic equation and singular LQR problem, Passivity and KYP lemma, Automatic control theory, Data estimation and filtering, nD-Systems, Graph and Metroid decompositions,  Optimal Control Behavioral Theory, Information encryption and coding error control, Power system stabilization, Analysis and simulation of very large scale (million node) circuits, Electromagnetic field computations for building circuits with better performance, Computational methods for systems theory, Embedded Control.


\item \href{https://www.ee.iitb.ac.in/~apel/} {\color{blue} \textbf{Applied Power Electronics Laboratory}}\\
    (\textit{Ground Floor, EE Building})\\
Professor In-charge - Prof. Vivek Agarwal\\
Non-Conventional Energy-Modelling the steady state and dynamic characteristics of the PV, Fuel
cell and wind energy sources, Power converter topologies for standalone and grid connected PV,
FC and wind systems, Maximum Power Point Tracking Schemes, Micro grid Power Qualit

Relevant Specializations - EE3\\
\item \href{https://www.ee.iitb.ac.in/web/research/labs/poweranser} {\color{blue} \textbf{Power Anser Lab}}\\
    (\textit{Ground Floor, EE Building})\\
Professor In-charge - Prof. Shreevardhan A. Soman\\
Relevant Specializations - EE3\\
ower Anser Labs was setup in April 2007 in collaboration with TCS. The association aims at
leveraging research outputs to full-fledged software product, primarily in the form of web ser-
vices, which can be used by power utilities.

\item \href{https://www.ee.iitb.ac.in/web/research/labs/powersys} {\color{blue} \textbf{Power Systems Laboratory}}\\
    (\textit{Ground Floor, EE Building})\\
Professor In-charge - Prof. Anil Kulkarni\\
Relevant Specializations - EE3\\
Fully Functional scaled down model of 2 Area 4 machine system used for demonstration of power
system dynamic phenomena i.e. frequency, voltage dynamics, effect of controllers like AVR,
FACTS and HVDC converters (variable impedance, voltage source and line-commutated con-
verters) with digital control.

\item  \textbf{Simulation Centre for Power Electronics and Power Systems}\\
    (\textit{Ground Floor,EE Building})\\
    Professor In-charge - Prof. Himanshu Bahirat\\
Relevant Specializations - EE3\\
Simulation centre for PEPS conducts workshops for people from different colleges and industries to train them on software that are crucial for power electronic and power system applications.This lab has different software like Matlab, PSCAD, PSIM, DigSilent and SaberRD. Apart from
this, this lab is also used by students, requiring intensive use of simulation software available here. This lab accommodates it's own software server. The models of power electronics and power system components are implemented on these software that are available in web site of simulation
centre.\\


\item \href{http://www.ncpre.iitb.ac.in/index.html} {\color{blue} \textbf{National Centre for Photovoltaic Research and Education}}\\
    (\textit{3rd Floor, Annex Building})\\
    Professor In-charge - Prof. B. G. Fernandes\\
Relevant Specializations - EE3\\
The National Centre for Photovoltaic Research and Education (NCPRE) at IIT Bombay was launched in 2010 with funding from the Ministry of New and Renewable Energy (MNRE) of the Government of India. The broad objectives of NCPRE are to provide R\&D and education support
for India’s ambitious 100 GW solar mission. The Centre is involved in both basic and applied research activities. These research activities include silicon solar cell fabrication and characterization, new materials for PV devices, energy storage and batteries for PV, development of power
electronic interfaces for solar PV systems, and module characterization and reliability.\\


\item \href{https://www.ee.iitb.ac.in/web/research/labs/fcl} {\color{blue} \textbf{Field Computation Laboratory}}\\
    (\textit{Ground Floor, EE Building})\\
    Professor In-charge - Prof. Shrikrishna A. Khaparde\\
Relevant Specializations - EE3\\
\\
Concurrent and visionary works related to generation, transmission and distribution of electrical energy are being carried out. Distributed generation, microgrids and renewable energy are the focus of research in the field of generation. Transmission expansion planning, ancillary service
management and power system stability are some of the studies undertaken by researchers in the lab.\\


\item \href{https://www.ee.iitb.ac.in/~idlab/Home\%20Page.html} {\color{blue} \textbf{Insulation Diagnostics Lab}}\\
    (\textit{2nd Floor, EE Building})\\
    Professor In-charge - Prof. S. V. Kulkarni\\
Relevant Specializations - EE3\\
\\
This laboratory has a 100 kV ac / 140 kV dc setup along with a partial discharge measurement setup. The academic activities include giving exposure to undergraduate and postgraduate students about various high voltage engineering aspects. Study of corona inception voltage for var-
ious electrode arrangements, demonstration of frequencies radiated by corona are some of experiments conducted as part of academic activities.\\


\item \href{https://www.ee.iitb.ac.in/web/research/labs/pedc} {\color{blue} \textbf{Power Electronics Design Centre (PEDC)}}\\
    (\textit{Ground Floor, EE Building})\\
    Professor In-charge - Prof. B. G. Fernandes\\
Relevant Specializations - EE3\\
\\
The Power Electronics Research Lab has equipment, instruments and computational facilities for research in power electronics, machines and drives, renewable energy field (photo-voltaic and wind) and their power conditioning and interconnection with the grid, power Converter topologies for grid connected and standalone PV and wind system, MPPT Schemes, Machine design, BLDC motor, Multi-Level Inverter, DC Microgrid.\\


\item \href{http://www.ee.iitb.ac.in/emlab/} {\color{blue} \textbf{Electrical Machines Lab}}\\
    (\textit{Ground Floor, EE Building})\\
    Professor In-charge - Prof. B. G. Fernandes\\
Relevant Specializations - EE3\\
\\
This undergraduate teaching laboratory has electric motor-generator sets, measuring instrumentation and power electronic drives. Students are exposed to concepts of machine characteristics and control through lab experiments.

\item \textbf{Stochastic Systems Lab}\\
    (\textit{4th Floor, GG Building})\\
    Professor In-charge - Prof. Nikhil Karamchandani, Prof. Sharayu Moharir, Prof. J K Nair\\
Relevant Specializations - EE1,EE5\\
\\
Description: In our lab, we focus on resource allocation for modern communication systems. Few are caching in distributed networks, Age of Information aware scheduling for IoT devices, coded caching.

\item \href{http://www.ee.iitb.ac.in/~zelerajesh/index.php/teaching-2-2-3-2/} {\color{blue} \textbf{Advanced Integrated Circuits and Systems Lab (aiCAS Lab)}}\\
    (\textit{1st Floor, EE Building})\\
Professor In-charge - Prof. Rajesh Zele \\
Relevant Specializations – EE1, EE5, EE6 (IC design, embedded system design, and MEMs based sensors/circuit development) \\
Research at aiCAS lab focuses in the areas of high performance Analog, RF, Mixed-signal IC design and Embedded Systems design. State of art hardware and software tools are used for developing integrated circuits using both RF-Analog and Digital design flow. The lab is being equipped with measurement instruments required for post-silicon validation of ICs. Integration of electronic components on PCB for verification and product (hardware-software) co-development is enabled with help of dedicated staff.
\end{enumerate}
\chapterimage{./pictures/faculty_old.jpg}
\chapter{Faculty Members}
\section{EE1: Communication and Signal Processing}
\begin{tabular}{p{4.5cm} p{9cm}}

\hline 
\hline 
Faculty & Research Interests \\ 
\hline
\hline 

\href{http://www.iitb.ac.in/en/employee/prof-abhay-karandikar}{\color{blue}{Prof. Abhay Karandikar }} & Control and Performance Modeling of Wireless Networks, Quality of Service and Resource Allocation in Wired/Wireless Networks, Next Generation Wireless Network Protocols (related to 802.16m, LTE-Advanced and 4G Standards), Co-operative Relay and Self Organizing Network,Carrier Ethernet and Mobile Backhaul, Rural Wireless Network. \\ 
\hline

\href{https://www.ee.iitb.ac.in/~asethi/}{\color{blue}{Prof. Amit Sethi }} & Computational pathology,Medical image analysis,Deep learning,Machine learning,Computer vision,Image processing,Signal processing. \\ 
\hline 

\href{https://www.ee.iitb.ac.in/~animesh/}{\color{blue}{Prof. Animesh Kumar }} & Signal processing, communication systems, applied statistics, SRAM reliability models. \\ 
\hline

\href{https://www.ee.iitb.ac.in/wiki/faculty/bikash}{\color{blue}{Prof. Bikash Kumar Dey }} & Information Theory, Coding Theory, Wireless communication. \\ 
\hline 

\href{https://www.ee.iitb.ac.in/~dmanju/}{\color{blue}{Prof. D. Manjunath }} & Computer and Communication Network Protocols, Systems and Algorithms Performance Modeling, Queueing Theory and Simulation, Stochastic
Systems. \\ 
\hline 

\href{https://www.ee.iitb.ac.in/wiki/faculty/gkumar}{\color{blue}{Prof. Girish Kumar }} & Microstrip antennas and arrays, Broadband antennas, Microwave integrated circuits, EMI/ EMC, RF communication circuits. \\ 
\hline 

\href{https://www.ee.iitb.ac.in/~gskasbekar/}{\color{blue}{Prof. Gaurav S. Kasbekar}} & Modeling, design and analysis of wireless networks, Game theoretic and economic aspects of spectrum allocation, Cognitive radio networks,
Lifetime and coverage problems in wireless sensor networks. \\ 
\hline 

\href{https://www.ee.iitb.ac.in/~jayakrishnan.nair/}{\color{blue}{Prof. Jaykrishnan U. Nair }} & Queueing theory, Communication networks, Heavy tails. \\ 
\hline

\href{https://www.ee.iitb.ac.in/wiki/faculty/jjohn}{\color{blue}{Prof. Joseph John}} & Analog and Digital Circuits, Optical Fiber Communications, Indoor Optical Wireless Systems, Modern Electronic Systems and Instrumentation. \\ 
\hline 

\href{https://www.ee.iitb.ac.in/~akumar/}{\color{blue}{Prof. Kumar Appaiah}} & Signal processing for communication, fibre optics, wireless communication. \\ 
\hline 

\end{tabular} 

\begin{tabular}{p{4.5cm} p{9cm}}
\hline 

\href{https://www.ee.iitb.ac.in/~manojg/}{\color{blue}{Prof. Manoj Gopalkrishnan}} & Algorithms in nature,Information processing in networks,Reaction networks,Neural networks,Evolution,Game theory,Deep learning,Information geometry,Thermodynamics of information,Quantum Information.\\ 
\hline 

\href{https://sites.google.com/site/nikhilkaram/}{\color{blue}{Prof. Nikhil Karamchandani }} & Information Theory, Networks, Communications,Distributed Computation, Cyber-Physical Systems. \\ 
\hline 

\href{https://www.ee.iitb.ac.in/wiki/faculty/chaporkar}{\color{blue}{Prof. Prasanna Chaporkar }} & Resource Allocation and scheduling in wired/wireless networks, Optimization and control of stochastic systems, Distributed systems and algorithms. \\ 
\hline

\href{https://www.ee.iitb.ac.in/wiki/faculty/prao}{\color{blue}{Prof. Preeti Rao }} & Speech and Audio Signal Processing, Music Information Retrieval. \\ \hline 

\href{https://www.ee.iitb.ac.in/~pcpandey/}{\color{blue}{Prof. Prem C. Pandey }} & Speech and Signal Processing, Biomedical Signal Processing and Instrumentation,Electronic Instrumentation, Embedded Electronic System Design. \\ 
\hline 

\href{https://www.ee.iitb.ac.in/web/faculty/homepage/rajbabu}{\color{blue}{Prof. Rajbabu Velmurugan }} & Statistical and digital signal processing, Signal processing system design, Particle filter applications, Target tracking systems. \\ 
\hline 

\href{https://www.ee.iitb.ac.in/~sarva/}{\color{blue}{Prof. Saravanan Vijayakumaran}} & Signal Processing for Communications, Parallel Simulation Algorithms. \\ 
\hline 


\href{https://www.ee.iitb.ac.in/wiki/faculty/shalabh}{\color{blue}{Prof. Shalabh Gupta}} & High-speed CMOS analog/RF/mm-wave integrated circuits and systems,Optical fiber communication systems, Microwave photonics / ultrafast data  conversion using photonics , Beam forming antenna systems, Signal processing for these systems. \\ 
\hline 

\href{https://www.ee.iitb.ac.in/wiki/faculty/bsraj}{\color{blue}{Prof. Sibi Raj B Pillai  }} & Fundamental Limits of Communication Systems, Information Theory and its applications, Compressed Sensing, Stochastic Modeling, Resource Allocation Problems, Interference Channels, Relaying and Broadcasting. \\ 
\hline 

\href{https://sites.google.com/site/sharayumoharir/}{\color{blue}{Prof. Sharayu Moharir }} & Modeling and the design of scalable resource allocation algorithms for large networks, including content delivery networks, communication
networks and crowd-sourcing. \\ 
\hline 

\href{https://www.ee.iitb.ac.in/wiki/faculty/merchant}{\color{blue}{Prof. Shabbir Merchant }} & Signal Processing, Adaptive Signal Processing. \\ 
\hline 

\href{https://www.ee.iitb.ac.in/~sc/}{\color{blue}{Prof. Subhasis Chaudhuri}} & Multimedia, Computer Vision, Image Processing, Pattern Recognition,Biomedical Signal Processing, Computational Haptics. \\ 
\hline 

\href{https://www.ee.iitb.ac.in/wiki/faculty/vmgadre}{\color{blue}{Prof. Vikram M. Gadre }} & Communication and signal processing, with emphasis on multiresolution
and multi-rate signal processing, especially wavelets and filter banks:
theory and applications. \\ 
\hline
\end{tabular} 

%\pagebreak

\section{EE2: Control and Computing}
\begin{tabular}{p{4.5cm} p{9cm}}

\hline 
\hline 
Faculty & Research Interests \\ 
\hline
\hline 

\href{https://www.ee.iitb.ac.in/wiki/faculty/dc}{\color{blue}{Prof. Debraj Chakraborty }} & Optimal Control, Linear Systems, Optimization, Differential Games,Game Theory \\ 
\hline 

\href{https://scholar.google.co.in/citations?user=9lWahYMAAAAJ&hl=en}{\color{blue}{Prof. Dwaipayan Mukherjee }} & Multi-agent Systems,Consensus,Formation Control,Control Theory and Robust Control \\ 
\hline 

\href{https://www.ee.iitb.ac.in/~debasattam/}{\color{blue}{Prof. Debasattam Pal}} & Distributed parameter systems, algebraic analysis, optimal control. \\ 
\hline 


\href{https://www.ee.iitb.ac.in/wiki/faculty/hp}{\color{blue}{Prof. Harish K. Pillai }} & Control theory, Systems theory, Multidimensional systems, Numerical
and computational methods, Coding theory, Optimization techniques, Electromagnetics \\ 
\hline
 

\end{tabular}


\begin{tabular}{p{4.5cm} p{9cm}}

\hline

\href{https://www.ee.iitb.ac.in/wiki/faculty/belur}{\color{blue}{Prof. Madhu N. Belur}} & Control theory, dissipative systems, graph theoretic methods, decentralized control, behavioral theory control, Fault diagnosis \\
\hline 

\href{https://www.ee.iitb.ac.in/wiki/faculty/chaporkar}{\color{blue}{Prof. Prasanna Chaporkar }} & Resource Allocation and scheduling in wired/wireless networks, Optimization and control of stochastic systems, Distributed systems and algorithms \\ 
\hline 

\href{https://www.ee.iitb.ac.in/~sc/}{\color{blue}{Prof. Subhasis Chaudhuri }} & Multimedia, Computer Vision, Image Processing, Pattern Recognition,Biomedical Signal Processing, Computational Haptics \\ 
\hline 

\href{https://www.ee.iitb.ac.in/wiki/faculty/svk}{\color{blue}{Prof. Shrikrishna V. Kulkarni}} & Transformers: Design, Analysis and Diagnostics, Electromagnetic and Coupled Field Computations, Power Engineering: Distributed Generation, High Voltage Engineering: Insulation Design/Diagnostics \\ 
\hline 


\href{https://www.ee.iitb.ac.in/web/faculty/homepage/borkar}{\color{blue}{Prof. Vivek Shripad Borkar }} & Stochastic Control, Learning Control Theory, Random Processes \\ 
\hline 

\href{https://www.ee.iitb.ac.in/wiki/faculty/vrs}{\color{blue}{Prof. Virendra R. Sule }} & Cryptology: Block and stream ciphers, efficient arithmetic for public key cryptography, algebraic cryptanalysis, Dynamical systems and feedback control theory \\ 
\hline

\end{tabular}


\section{EE3: Power Electronics and Power Systems}
\begin{tabular}{p{4.5cm} p{9cm}}

\hline 
\hline 
Faculty & Research Interests \\ 
\hline 
\hline 


\href{https://www.ee.iitb.ac.in/wiki/faculty/anil}{\color{blue}{Prof. Anil Kulkarni }} & Power System Dynamics and Control, Application of Power Electronics to Power Systems, Flexible AC Transmission Systems \\ 
\hline 


\href{https://www.ee.iitb.ac.in/wiki/faculty/ashukla}{\color{blue}{Prof. Anshuman Shukla }} & Multilevel converters and Modulation and control of power electronic converters, Power electronics applications in power systems (FACTS,HVDC, custom power devices, etc.), Renewable energies and Energy storage, Control of electric drives, Hybrid and solid-state circuit breakers and current limiters \\ 
\hline 


\href{https://www.ee.iitb.ac.in/web/faculty/homepage/anu}{\color{blue}{Prof. Anupama Kowli }}  & Power System Planning, Operations and Control, Electricity Markets and Economics of Electric Power Grids, Demand-side Management, Demand
Response and Flexible Loads, Smart Grids and its Enabling Technologies and Mechanisms, Policy and Regulation for Electric Power Grids. \\ 
\hline

\href{https://www.ee.iitb.ac.in/wiki/faculty/bgf}{\color{blue}{Prof. Baylon G. Fernandes}} & Inverter topologies for VAR compensation, Power electronic interface for non-conventional energy sources, Permanent magnet machines for
wind power generation, Switched reluctance machines for electric vehicle application \\ 
\hline 

 
\href{https://www.ee.iitb.ac.in/web/faculty/homepage/hjbahirat}{\color{blue}{Prof. Himanshu J. Bahirat }} & Renewable Energy Sources; Grid Integration of Renewable Energy; Offshore Wind Energy; Transients in Power Systems; DC Power Systems;
DC Wind Farms; Multi-terminal DC Networks; Circuit Breakers; Power Electronics. \\ 
\hline 

\href{https://www.ee.iitb.ac.in/wiki/faculty/kishore}{\color{blue}{Prof. Kishore Chatterjee}} & Utility friendly converter topologies, Power Factor Correction techniques,STATCOM, Switched Mode Rectifiers, Electronic Ballast \\ 
\hline 


\href{https://www.ee.iitb.ac.in/~mukul/}{\color{blue}{Prof. Mukul C. Chandorkar }} & Power Electronics, Power quality, Static Compensation, Motor Drives \\ 
\hline 

\href{https://www.ee.iitb.ac.in/~pcpandey/}{\color{blue}{Prof. Prem C. Pandey}} & Speech and Signal Processing, Biomedical Signal Processing and Instrumentation,Electronic Instrumentation, Embedded Electronic System Design \\ 
\hline 
\end{tabular}


\begin{tabular}{p{4.5cm} p{9cm}}


\href{https://www.ee.iitb.ac.in/wiki/faculty/soman}{\color{blue}{Prof. Shreevardhan A. Soman}}  & Power system analysis, computation and economics, Power system protection \\ 
\hline 

\href{https://www.ee.iitb.ac.in/wiki/faculty/sak}{\color{blue}{Prof. Shrikrishna A. Khaparde }} & Deregulation in Power Industry: optimal bidding, and congestion management,Object Oriented Power System Analysis, Controlled series compensation using SSSC, Harmonic Distortion in Distribution systems, Design and Operation of small tidal power plant, Modeling and Design of transformer \\ 
\hline 

\href{https://www.ee.iitb.ac.in/wiki/faculty/svk}{\color{blue}{Prof. Shrikrishna V. Kulkarni }} & Transformers: Design, Analysis and Diagnostics, Electromagnetic and Coupled Field Computations, Power Engineering: Distributed Generation, High Voltage Engineering: Insulation Design/Diagnostics \\ 
\hline 

\href{https://www.ee.iitb.ac.in/~agarwal/}{\color{blue}{Prof. Vivek Agarwal }} & Power conversion: New converter topologies, High frequency link
power conversion, ZCS-ZVS configurations, Switched Capacitor DCDC converters, Power quality issues: Power factor correction techniques, Static VAR compensation, Active filters \\ 
\hline 

\end{tabular} 



\section{EE5: Electronic Systems}
\begin{tabular}{p{4.5cm}p{9cm}}

\hline 
\hline 
Faculty & Research Interests \\ 
\hline
\hline


 \href{https://www.ee.iitb.ac.in/wiki/faculty/dinesh}{\color{blue}{Prof. Dinesh K. Sharma}} & MOS device modeling VLSI design and technology. Microelectronics - technology and device characterization mixed signal design \\ 
 \hline 

\href{https://www.ee.iitb.ac.in/wiki/faculty/jjohn}{\color{blue}{Prof. Joseph John}} & Analog and Digital Circuits, Optical Fiber Communications, Indoor Optical Wireless Systems, Modern Electronic Systems and Instrumentation \\ 
\hline 

\href{https://www.ee.iitb.ac.in/~madhav/}{\color{blue}{Prof. Madhav P. Desai}} & VLSI Circuits and Systems, VLSI design and design automation,Graph theory and combinatorics.\\ 
\hline 

 \href{https://www.ee.iitb.ac.in/wiki/faculty/mshojaei}{\color{blue}{Prof. Maryam Shojaei Baghini}} & Analog/Mixed-signal VLSI design and test (SoC, LV, LP, LE, Bio-medical/Biosensors, Bio-inspired circuits and systems, I/O, highly precise circuits and systems, instrumentation, energy harvesting and many more applications), Specific technologies and performance-optimized Analog/mixed-signal/RF circuits and systems for healthcare applications.\\ 
\hline 

\href{https://www.ee.iitb.ac.in/~pcpandey/}{\color{blue}{Prof. Prem C. Pandey}} &  Speech and Signal Processing, Biomedical Signal Processing and Instrumentation, Electronic Instrumentation, Embedded Electronic System Design\\ 
\hline 

 \href{https://www.ee.iitb.ac.in/wiki/faculty/prao}{\color{blue}{Prof. Preeti Rao}} &  Speech and Audio Signal Processing, Music Information Retrieval \\ 
 \hline 

 \href{https://www.ee.iitb.ac.in/web/people/faculty/home/rajeshzele}{\color{blue}{Prof. Rajesh H. Zele }} & RF, Analog and Mixed-Signal Circuits for Communication Applications. \\ 
\hline 

\href{https://www.ee.iitb.ac.in/wiki/faculty/patkar}{\color{blue}{Prof. Sachin Patkar}} & Combinatorial optimization Matroid Theory Submodular Functions
Linear/Integer programming Network Flows High Performance Computing FPGA-based accelerated computing GPU based acceleration High Performance Circuit Simulation Algorithms Design and Analysis \\ 
 \hline 

\href{https://www.ee.iitb.ac.in/wiki/faculty/shalabh}{\color{blue}{Prof. Shalabh Gupta }} & High-speed CMOS analog/RF/mm-wave integrated circuits and systems,Optical fiber communication systems, Microwave photonics / ultrafast data  conversion using photonics , Beam forming antenna systems, Signal processing for these systems \\ 
\hline 

\href{https://www.ee.iitb.ac.in/wiki/faculty/merchant}{\color{blue}{Prof. Shabbir Merchant }} & Signal Processing, Adaptive Signal Processing\\ 
\hline 

\end{tabular}


\begin{tabular}{p{4.5cm} p{9cm}}

\hline
\href{https://www.ee.iitb.ac.in/web/people/faculty/home/sdgupta}{\color{blue}{Prof. Siddhartha P. Duttagupta }} & Microelectronics, Micro/Nano Sensor Technology Optimization and Application, Sensor Integrated Electronic Circuits and Systems- Design \\ 
\hline 

 \href{http://www.ee.iitb.ac.in/~stallur/index.php}{\color{blue}{Prof. Prof. Siddharth Tallur}} &  RF MEMS, Photonics, Opto-Mechanics, Micro- and Nano-fabrication, Sensor Systems.\\ 
\hline 

\href{https://www.ee.iitb.ac.in/wiki/faculty/vmgadre}{\color{blue}{Prof. Vikram M. Gadre}} & Communication and signal processing, with emphasis on multiresolution and multi-rate signal processing, especially wavelets and filter banks: theory and applications.\\ 
\hline 

\end{tabular}

\bigbreak

\section{EE6: Integrated Circuits}
\begin{tabular}{p{4.5cm} p{9cm}}
\hline 
\hline
Faculty  & Research Interests \\ 
\hline
\hline

\href{https://www.ee.iitb.ac.in/wiki/faculty/dinesh}{\color{blue}{Prof. Dinesh K. Sharma}} & MOS device modeling VLSI design and technology. Microelectronics - technology and device characterization mixed signal design \\
\hline

\href{https://www.ee.iitb.ac.in/~madhav/}{\color{blue}{Prof. Madhav P. Desai}} & VLSI Circuits and Systems, VLSI design and design automation,Graph theory and combinatorics.\\ 
\hline

\href{https://www.ee.iitb.ac.in/wiki/faculty/mshojaei}{\color{blue}{Prof. Maryam Shojaei Baghini}} & Analog/Mixed-signal VLSI design and test (SoC, LV, LP, LE, Bio-
medical/Biosensors, Bio-inspired circuits and systems, I/O, highly precise circuits and systems, instrumentation, energy harvesting and many more applications), Specific technologies and performance-optimized Analog/mixed-signal/RF circuits and systems for healthcare applications.\\ 
\hline

\href{https://www.ee.iitb.ac.in/web/people/faculty/home/jayanta}{\color{blue}{Prof. Jayanta Mukherjee}} & RF VLSI Design, Antenna Design, Biomedical IC Design, Testing and Characterization, Analog VLSI, Digital VLSI, Noise Modeling.\\
\hline

\href{https://www.ee.iitb.ac.in/web/people/faculty/home/pramodm}{\color{blue}{Prof. Pramod Murali }} & CMOS Biosensors, Sensors for Autonomous Drones, Robotic Prostheses, Artificial Intelligence (AI) at Sensor Nodes.\\
\hline

\href{https://www.ee.iitb.ac.in/web/people/faculty/home/rajeshzele}{\color{blue}{Prof. Rajesh H. Zele }} & RF, Analog and Mixed-Signal Circuits for Communication Applications. \\
\hline

\href{https://www.ee.iitb.ac.in/wiki/faculty/shalabh}{\color{blue}{Prof. Shalabh Gupta }} & High-speed CMOS analog/RF/mm-wave integrated circuits and systems,Optical fiber communication systems, Microwave photonics / ultrafast data  conversion using photonics , Beam forming antenna systems, Signal processing for these systems.\\ 
\hline

\href{https://www.ee.iitb.ac.in/~viren/}{\color{blue}{Prof. Virendra Singh}} & Computer Architecture Processor architecture and micro-architecture
VLSI Testing Fault-tolerant computing Robust design and architectures Self-healing system design SoC/NoC design and test Post Silicon Debug High level synthesis Formal verification. \\ 
\hline
\end{tabular}

\section{EE7: Solid State Devices}
\begin{tabular}{p{3.5cm} p{9cm}}

\hline 
\hline
Faculty  & Research Interests \\ 
\hline

\href{https://www.ee.iitb.ac.in/web/people/faculty/home/anilkg}{\color{blue}{Prof. Anil Kottantharayil}} & CMOS device physics, design and modeling, Materials for advanced CMOS devices CMOS device physics, design and modeling, Materials for advanced CMOS devices, Electrical characterization.\\
\hline


\href{https://www.ee.iitb.ac.in/web/people/faculty/home/laha}{\color{blue}{Prof. Apurba Laha}} & Oxide based electronics: Epitaxial rare earth oxides (high-K dielectrics) on Si, Ge and SiGe substrates for gate dielectric application in next generation CMOS, oxide hetero-structures.\\
\hline

\end{tabular}


\begin{tabular}{p{4.5cm} p{9cm}}
\hline
\href{https://www.ee.iitb.ac.in/web/people/faculty/home/ashwin}{\color{blue}{Prof. Ashwin A. Tulapurkar}} & Spintronics, Physics of nano-devices, Spin-current induced magnetization switching, RF properties of spintronic devices, Noise, Thermoelectric effects.\\
\hline

\href{https://www.ee.iitb.ac.in/web/people/faculty/home/bm}{\color{blue}{Prof. Bhaskaran Muralidharan}} & Computational nanolectronics, spintronics , nano scale energy conversion.\\
\hline

\href{https://www.ee.iitb.ac.in/wiki/faculty/dinesh}{\color{blue}{Prof. Dinesh K. Sharma}} & MOS device modeling VLSI design and technology. Microelectronics - technology and device characterization mixed signal design \\
\hline

\href{https://www.ee.iitb.ac.in/web/people/faculty/home/dsaha}{\color{blue}{Prof. Dipankar Saha}} & Microelectronics New Device Physics Semiconductor Spintronics Spin injection, transport and detection in III-V systems Device Reliability.\\
\hline

\href{https://www.ee.iitb.ac.in/web/people/faculty/home/dm}{\color{blue}{Prof. Dwaipayan Mukherjee}} & Control Theory, Robust Control, Multi-agent Systems, Consensus protocols, Formation control.\\
\hline


\href{https://www.ee.iitb.ac.in/web/people/faculty/home/kasturis}{\color{blue}{Prof. Kasturi Saha}} & Nano photonics, Quantum optics, Magnetometry, Colour centres in diamond.\\
\hline

\href{https://www.ee.iitb.ac.in/web/people/faculty/home/mbpatil}{\color{blue}{Prof. Mahesh B. Patil}} & Circuit simulation Semiconductor device modeling and simulation Real-time simulation of power electronic circuits and systems.\\
\hline

\href{https://www.ee.iitb.ac.in/~rrao/}{\color{blue}{Prof. Valipe Ramgopal Rao}} & Nanoelectronics Technology Aware Design Challenges with Emerging Technologies (Process-Device-Circuit Interactions with Multigate MOSFETs, Polymer Transistors, Molecular Electronics, etc.) CMOS Reliability Bio-MEMS.\\
\hline

\href{https://www.ee.iitb.ac.in/web/people/faculty/home/naren}{\color{blue}{Prof. Narendra Sudhir Shiradkar}} & Reliability of solar cells, modules and systems, PV Performance Monitoring, Design for reliability, , Techniques for PV Power Plant Inspection and Monitoring, Power Electronic Device reliability.\\
\hline

\href{https://www.ee.iitb.ac.in/web/people/faculty/home/prnair}{\color{blue}{Prof. Pradeep R. Nair}} & Nanoscale devices for energy and healthcare applications., Semiconductor device physics and reliability., Micro Electro Mechanical systems (MEMS).\\
\hline

\href{https://www.ee.iitb.ac.in/web/people/faculty/home/slodha}{\color{blue}{Prof. Saurabh Lodha}} & CMOS process integration and device physics, Materials and processes for advanced CMOS devices, Metal-semiconductor interfaces.\\
\hline

\href{https://www.ee.iitb.ac.in/web/people/faculty/home/souvik}{\color{blue}{Prof. Souvik Mahapatra}} & Electrical characterization, modeling and simulation of micro/nano electronic devices, NBTI/PBTI and Hot carrier degradation in MOSFETs, High-k gate dielectrics, Advanced CMOS device reliability.\\
\hline

\href{https://www.ee.iitb.ac.in/web/people/faculty/home/sdgupta}{\color{blue}{Prof. Siddhartha P. Duttagupta}} & Microelectronics, Micro/Nano Sensor Technology Optimization and Application, Sensor Integrated Electronic Circuits and Systems- Design, Development and Deployment.
\\
\hline

\href{https://www.ee.iitb.ac.in/web/people/faculty/home/sganguly}{\color{blue}{Prof. Swaroop Ganguly}} & Physics and technology of nanoscale devices, Spin-based devices and circuits, Energy-conversion devices.\\
\hline

\href{https://www.ee.iitb.ac.in/web/people/faculty/home/subho}{\color{blue}{Prof. Subhananda Chakrabarti}} & III-V Compound semiconductor materials growth and characterization, Optoelectronic Devices of interest includes quantum dot photo detectors and Solar cells, III-V device integration on germanium.\\ 
\hline

\href{https://www.ee.iitb.ac.in/web/people/faculty/home/udayan}{\color{blue}{Prof. Udayan Ganguly}} & Flash Memory materials and device optimization, Nanocrystal, nanotubes, nanowires growth and device characterization, Novel materials/process integration in CMOS platform.\\
\hline

\end{tabular}

%----------------------------------
%chapter 5 ends
%---------------------------------------



\chapterimage{./pictures/dep_act_bg.jpg}
\chapter{Department Activities}
\section {Student's Reading Group (SRG):}
The SS Students’ Reading Group (SRG) was started in 2015 as an interactive peer review basedplatf orm to share knowledge and research issues from various domains of Electrical Engineering. The 5 specializations of EE department are divided into 4 clusters and researchers from every cluster present and discuss their work with their peer in the talks that are conducted throughout the semester. The sessions are entirely student run, giving the speakers a unique opportunity to present their ideas freely and receive reviews from the students alone.\newline
New phase of SRG begins every semester. A 5 minute research challenge (5MR) is conducted in each phase where the participants get 5 minute to put forward their research ideas. A panel of faculty members decides the best speakers who are then awarded. The 5MR challenge is aimed
at improving the technical communication skills of the student.
\section {Department Academic Assistance Program (DAAP):}
Department Academic Assistance Program (DAAP) is a helping hand to the students who are facing challenges in their academics. A student can face academic related issues due to various reasons, like unacquaintance with highly competitive environment at IIT, managing the academic
workload and assignment timelines, unclear about field of interest etc. We believe that a little help from an experienced friend can make a difference. We provide one to one assistance through regular meetings. Assistants are student who have already excelled in the concerned field/course. They can complement classroom learning by providing techniques to manage course content effectively and sharing the resources. Since the inception of this program we have observed an improvement in the performance of more than 15 students. DAAP is creating a better environment for knowledge sharing and learning with joy!
\newline
\section {IEEE Students Chapter:}
The IIT Bombay IEEE Students’ Chapter is a student body that strives to promote excellence in various fields of electrical engineering. The organization is involved in organizing talks, workshops and other activities to help students obtain new skills in their fields of interest.\newline
\section {Bridge Course:}
The Bridge Course, starting from 2016 is an initiative by the department to help new students of the M. Tech. and PhD programs which makes them comfortable with their coursework. The Bridge course focuses on revising essential prerequisites and developing analytical skill. It is felt that these two essential skills will help students to work on their courses more effectively.\newline
\section {AAVRITI (Department Techfest):}
It is the annual research and technological festival of Electrical Engineering department of IIT Bombay. AAVRITI, with the motto of promoting technology, creativity, intelligence and sheer innovation, and to enthrall the magical impact of electronics that it has on the human civilization over the years aims to bridge the gap between industry \& academia by encouraging exchange of ideas and providing opportunities for technical interactions.\newline
It includes workshops on most advanced and buzzing technologies, lecture series by experts of particular fields as well as competitions and hackathons. This gives students a chance to expand their horizons and learn beyond engineering syllabus.\newline
Other department activities include Teacher’s day, Sports day, Department trip, Convocation ceremony, Valedictory function and so on.

\chapterimage{./pictures/logos.jpg} % Chapter heading image
\chapter{Department Representatives}

\section{Department placement coordinators (2018-2019) }\index{Department placement coordinators }\par
\begin{center}
	\begin{tabular}{cccc}
		Pulkit Rajgadiya  &  Mudit Vijayvargia   &	 Shashank Joshi & Nayan Nitnaware \\  
		 \textit{183070001@iitb.ac.in}   &  \textit{183070035@iitb.ac.in}   &  \textit{183070064@iitb.ac.in}   &\textit{183070019@iitb.ac.in}    \\
		     9983481233     &   9034118076  &       8057241101  & 9096140415 \\

	\end{tabular}
\end{center}
\sectionlinetwo{magenta}{85}


\section{Company coordinators}\index{Company coordinators}
\begin{center}
	\begin{tabular}{ccc}
		Sourabh Suri    & 	 Shikha Das    & Sourabh Patil  \\
		\textit{sourabh94suri@gmail.com}       &      \textit{shikha.sd1996@gmail.com}   	& \textit{sourabh.p.iitb@gmail.com}\\
		    8826533683   &      9644147411 & 8806070395         \\
	\end{tabular}
\end{center}
\sectionlinetwo{magenta}{85}
 %\newpage

\section{EE Student Assosciation}\index{EE Student Assoscition}
Posts relevant to a fresher Mtech student are mentioned here. For more information about EESA visit \href{https://www.ee.iitb.ac.in/course/~eesa}{\color{blue}{EESA website}}
\begin{center}
	\begin{tabular}{ccc}
		  \textbf{SRG Overall-Cordinator    }      &   \textbf{General Secretary   }   &   \textbf{Cultural Secretary   }\\
		          Shashank Kurm                &          Abhijeet Anand           &      Prasun Sukai         \\
		     \textit{oc\_srg@ee.iitb.ac.in}  &  \textit{dgsec@ee.iitb.ac.in}    &  \textit{pgcultsec@ee.iitb.ac.in}  \\
		          9582499323                          &           829147642               &        8895109637                \\
	\end{tabular}
\end{center}

%\newpage
\bigbreak
\begin{center}
	\begin{tabular}{cc}
		  \textbf{Sports Secretary (Boys)}      &     \textbf{Sports Secretary (Girls)}\\
		         Prashant Sharma                &               BVS Anusha \\
		\textit{183079037@iitb.ac.in} &\textit{bvsanusha@ee.iitb.ac.in} \\
		            9911155297                  &                    9869607702 \\
	\end{tabular}
\end{center}
\sectionlinetwo{magenta}{88}

\newpage
\section{EE ISCP Team}\index{EE ISCP Team}

%\textbf {Institute Student Companion Programme (2019-20)}\\
\medskip
\begin{center}
\bigbreak
\large{\textbf{Department Coordinators}}
\bigbreak
	\begin{tabular}{cc}
		\photo{1cm}{35mm}{./pictures/sunny.jpeg}
		& \photo{1cm}{35mm}{./iscp/ram.jpg} \\
		 \textbf{Sunny Mehta}
		&\textbf{Ramkrishna}\\
		\textit{sunnymehta78669@gmail.com}
		&\textit{ramporicha@ee.iitb.ac.in}\\
		 +919571648532 & +919677165155 
	\end{tabular}
\end{center}
\bigbreak
\bigbreak
\begin{center}
\bigbreak
\large{\textbf{Student Companions}}

\bigbreak
\bigbreak
\bigbreak
	\begin{tabular}{ccc}
		\photo{1cm}{35mm}{./iscp/sitaram.jpg}
		& \photo{1cm}{32mm}{./iscp/indrani.jpg}
		& \photo{1cm}{35mm}{./iscp/tarun.jpg}\\
		  \textbf{G Nagasitaram}
		& \textbf{Indrani Mukherjee}
		& \textbf{Tarun S} \\
		\textit{sitaram@ee.iitb.ac.in}
		&\textit{mukherjee.indrani22@gmail.com}
		&\textit{tarunsathesh@gmail.com}\\
		\photo{1cm}{31mm}{./iscp/aswin.jpg}
		& \photo{1cm}{35mm}{./iscp/sourabh.jpg}
		& \photo{1cm}{25mm}{./iscp/pankaj.jpg}\\
		   \textbf{Aswin Ajayan}
		&  \textbf{Patil  Sourabh}
		&  \textbf{Pankaj Singh}\\
		\textit{aswin@ee.iitb.ac.in}
		&\textit{sourabh.p.iitb@gmail.com}
		&\textit{contact.pankaj.singh7@gmail.com}\\
		\photo{1cm}{35mm}{./iscp/amey.jpg}
		& \photo{1cm}{30mm}{./iscp/raman.jpg}
		& \photo{1cm}{35mm}{./iscp/yaswanth.jpg}\\
		  \textbf{Chindarkar Amey}
		& \textbf{Raman Thukral} 
		& \textbf{Yaswanth Chebrolu}\\
		\textit{amey2994@gmail.com}
		&\textit{ramanthukral111@gmail.com}
		&\textit{yaswanth.chebrolu@gmail.com}\\
\end{tabular}
\end{center}

\newpage
\begin{center}
	\begin{tabular}{ccc}

		\photo{1cm}{26mm}{./iscp/risabh.jpg}
		& \photo{1cm}{35mm}{./iscp/jahnavi.jpg}
		& \photo{1cm}{35mm}{./iscp/sabitha.jpg}\\
		  \textbf{Risabh Chana}
		& \textbf{N Jahnavi}
		& \textbf{Sabitha Joseph}	\\	 
		\textit{rishabhchana844@gmail.com}
		&\textit{jahnavireddy924@gmail.com}
		&\textit{sabi.joseph3@gmail.com}\\


		\photo{1cm}{35mm}{./iscp/ashvini.jpg}
		& \photo{1cm}{27mm}{./iscp/pulkit.jpg}
		& \photo{1cm}{30mm}{./iscp/nijil.jpg}\\
		  \textbf{Ashvini Kumar}
		& \textbf{Pulkit Jain}
		& \textbf{Nijil George}	\\	 
		\textit{sharmaashvinikumar8@gmail.com}
		&\textit{jainpulkit54@gmail.com}
		&\textit{nijiliitb@gmail.com}\\
		\photo{1cm}{35mm}{./iscp/sancho.jpg}
		& \photo{1cm}{35mm}{./iscp/samiksha.jpg}
		& \photo{1cm}{35mm}{./iscp/rajat.jpg}\\
		  \textbf{Sancho Simmy Louis}
		& \textbf{Samiksha Baid}
		& \textbf{Rajat Chiddarwar}\\	 
		\textit{sancho@iitb.ac.in}
		&\textit{sami71095.sb@gmail.com}
		&\textit{rajatchiddarwar123@gmail.com}\\

	\end{tabular}
\end{center}

\bigbreak
\bigbreak

\section{EE Office Contact Persons}
%\medskip

\begin{center}
	\begin{tabular}{cccc}
\textbf{Madhummati Shetty  }    &   \textbf{Santosh S. Kharat   }   &   \textbf{Tanvi D. Shelatkar   }  &   \textbf{Soni Dhochak } \\           
   madhus@ee.iitb.ac.in         &     santoshk@ee.iitb.ac.in        &       tanvi@e.iitb.ac.in          &      soni@ee.iitb.ac.in  \\
      (022)2576-7401            &       (022)2576-7402              &         (022)2576-7401            &         (022)2576-7402   \\
	\end{tabular}
\end{center}

\bigbreak
\bigbreak
\chapterimage{./pictures/useful-links.jpg}
\chapter{Links you should check out}
\section{Links you should check out}
\begin{itemize}
	\item \href{https://www.ee.iitb.ac.in}{\color{blue}{Department Website}}
	\item Explore more at \href{http://www.iitb.ac.in/en/activities/student-clubs}{\color{blue}{Clubs at IITB}}
	\item \href{https://insti.app/feed}{\color{blue}{Insti App}} Know your Campus
	\item \href{https://gymkhana.iitb.ac.in/~scp/scp/index.html}{\color{blue}{Institute Students Companion Program (ISCP)}}
	%\item \href{http://ftp.iitb.ac.in/}{\color{blue}{IIT Bombay FTP Server}} Get daiy builds of various linux distributions here and more
	\item \href{http://ftp.iitb.ac.in/}{\color{blue}{IIT Bombay FTP Server}} A common repository of  general useful stuff for IITB community like daily builds of various linux distributions and more.
	\item \href{https://www.cc.iitb.ac.in/}{\color{blue}{Computer Center}} To get licenses for various advanced and special purpose softwares such as Ansys, Matlab, Mathematica, Maple, Techplot, etc. for all the campus users
	\item \href{http://www.ee.iitb.ac.in/~pclab/index.php/}{\color{blue}{PC Lab}} A number of computers with Windows® as well as Linux are available for students to use
	
\end{itemize}


%\textit{Wish you all the best, ISCP Team}

\end{document}
